\section{Results}\label{sec:results}

The results are reported in the order of the four questions of Section~\ref{sec:intro}, each answered by the audit steps assigned to it in Section~\ref{sec:intro}, followed by the attribution summaries of the rule-only variant and by the robustness checks. Unless stated otherwise, metrics are means over the 37 evaluable folds at the transferred-threshold operating point, ROC-AUC is threshold-invariant, and paired tests are Wilcoxon signed-rank tests over folds with Holm correction within each family (Section~\ref{sec:eval}).

\subsection{The Gaze Inputs Account for the Accuracy}\label{res:rq1}

\paragraph{The audited decoder performs comparably to the best tested baselines.} Under LOSOCV the decoder reached an accuracy of $78.2\pm16.6\%$, a balanced accuracy of $0.75\pm0.18$, an ROC-AUC of $0.88\pm0.17$ and an MCC of $0.49\pm0.36$ over the 37 held-out participants, with subject-level bootstrap 95\% confidence intervals of $[0.82, 0.93]$ for ROC-AUC, $[0.69, 0.81]$ for balanced accuracy and $[0.37, 0.60]$ for MCC (Fig.~\ref{fig5}). At the calibrated operating point balanced accuracy was 0.72 and MCC 0.45, with ROC-AUC unchanged; the difference is a threshold effect. A within-fold label permutation (5{,}000 shuffles) gave $p=2\times10^{-4}$ for mean per-fold ROC-AUC, balanced accuracy and MCC alike, so the association between the held-out predictions and the labels is not attributable to chance. Three trainings of the same configuration with the same seed gave balanced accuracy 0.72--0.75 and ROC-AUC 0.86--0.88, which is the run-to-run spread reported beside the decoder's own variants below. The pooled expected calibration error was 0.14 and the Brier score 0.18.

\begin{figure*}[!htbp]
\centering
\includegraphics[width=\textwidth]{fig3_losocv_results_r2.pdf}
\caption{LOSOCV performance (37 participants). (A) Pooled ROC curves of the audited decoder, the tuned gaze-only ET-LSTM and the gaze-free variant; the legend gives pooled AUC, and the per-fold means are those of the text. (B) Per-fold distributions of the decoder's metrics at the calibrated operating point (diamonds: means). (C) Mean per-fold ROC-AUC ($\pm$ SD) of the label's gaze and EEG terms under a linear probe (Supplementary Table~S15) beside the decoder, the tuned ET-LSTM, the gaze-free variant and the best tuned EEG encoder.}
\label{fig5}
\end{figure*}

\paragraph{The index is recoverable from its gaze terms.} A logistic probe under LOSOCV recovers the label from its five gaze terms alone at a mean per-fold ROC-AUC of 0.92 (pooled 0.94), from its five EEG terms alone at 0.67 (pooled 0.61), and from all ten terms at 1.00 (Supplementary Table~S15; audit step~1). These probes are the references for the two modalities. A score at or below 0.92 from a model that sees gaze, the decoder's 0.88 included, has not been shown to use information beyond the rule's gaze terms; a score at or below 0.67 from a model that sees only EEG has not been shown to use information beyond the rule's EEG terms. The comparison orders performance and does not identify the information a model uses: a model can score below a probe while using different information, or above it by using the same terms non-linearly.

\paragraph{No model of the principal comparison that sees gaze exceeded the gaze-only LSTM.} Table~\ref{tab:gaze} places every model that receives gaze on one footing (audit step~3). The three tuned eye-tracking encoders reach balanced accuracy 0.75--0.77 and ROC-AUC 0.88--0.89, the six fusion architectures 0.73--0.77 and 0.83--0.88, and the decoder 0.75 and 0.88. The two fusions built from standard parts sit at the lower end of that range: EEGNet+ET-LSTM reaches 0.73 and 0.81 and ShallowConvNet+ET-LSTM 0.72 and 0.83, with MCC 0.46 and 0.43. Under Holm-corrected paired tests over the 37 folds no difference between the decoder and any of the nine eye-tracking or fusion baselines survives correction on balanced accuracy, MCC or ROC-AUC (all $p_{\mathrm{Holm}}=1.0$; smallest uncorrected $p=0.12$, late fusion on ROC-AUC; $|r_{\mathrm{rb}}|\le0.39$; 15 to 21 of the 37 folds tied on ROC-AUC; Supplementary Table~S10). Against the tuned ET-LSTM the decoder differs by $-0.011$ in ROC-AUC ($p_{\mathrm{Holm}}=1.0$, $r_{\mathrm{rb}}=-0.05$), $-0.009$ in balanced accuracy and $-0.032$ in MCC (both $p_{\mathrm{Holm}}=1.0$; Supplementary Table~S14). Against the two fusions of standard parts, with Holm correction over the five new comparisons per metric, the decoder differs by $+0.016$ and $+0.029$ in balanced accuracy and by $+0.026$ and $+0.060$ in MCC (all $p_{\mathrm{Holm}}=1.0$), and in ROC-AUC by $+0.067$ against EEGNet+ET-LSTM (95\% CI $[+0.009, +0.121]$; raw $p=0.021$, $p_{\mathrm{Holm}}=0.082$; $r_{\mathrm{rb}}=0.53$; 20 wins, 12 ties, 5 losses) and $+0.045$ against ShallowConvNet+ET-LSTM ($[-0.018, +0.115]$; $p_{\mathrm{Holm}}=0.45$). The fusions rank epochs slightly less well than the gaze model they contain: adding EEGNet to the ET-LSTM lowers ROC-AUC by 0.078 ($[0.037, 0.124]$; $p_{\mathrm{Holm}}=0.006$; $r_{\mathrm{rb}}=0.77$; 19 wins, 14 ties, 4 losses) and adding ShallowConvNet by 0.056 ($[0.010, 0.110]$; $p_{\mathrm{Holm}}=0.17$), with no detectable change in balanced accuracy or MCC ($p_{\mathrm{Holm}}\ge0.40$); under decision-level averaging, an EEG encoder that is itself near chance on the index dilutes the gaze ranking. In the multimodal comparison, therefore, the eye-tracking LSTM's 0.89 was exceeded neither by the decoder nor by any fusion, the two fusions of standard parts ranked below it, and the corrected comparisons detected no advantage of the decoder over the gaze-only reference; the non-significant differences do not establish equivalence. Calibration does not separate them either: the decoder's pooled ECE of 0.14 and Brier score of 0.18 lie within the 0.06--0.15 and 0.15--0.19 of the gaze-based baselines, most of which are as well or better calibrated. Tuning decides this comparison: against an ET-LSTM trained under the common fixed budget (Supplementary Table~S8) the decoder leads numerically, a lead that disappears once the encoder is tuned per fold (0.89).

\begin{table*}[!htbp]
\caption{Every model that sees gaze, under LOSOCV with baselines tuned per fold (37 folds, transferred-threshold operating point; ECE and Brier pooled over the 347 held-out epochs). Label-coupled, since five label terms are gaze statistics. Paired tests against the decoder: Supplementary Table~S10 for the eye-tracking and fusion baselines and Section~\ref{res:rq1} for the two standard-encoder fusions (none survives Holm correction).}\label{tab:gaze}
\small
\begin{tabular*}{\linewidth}{@{\extracolsep{\fill}}lccccc@{}}
\toprule
Model & BalAcc & ROC-AUC & MCC & ECE & Brier \\
\midrule
\multicolumn{6}{@{}l}{\textit{Eye-tracking encoders (gaze only)}} \\
ET-Transformer & 0.77$\pm$0.19 & 0.88$\pm$0.17 & 0.54$\pm$0.39 & 0.07 & 0.16 \\
ET-LSTM & 0.76$\pm$0.19 & 0.89$\pm$0.15 & 0.52$\pm$0.37 & 0.10 & 0.16 \\
ET-GRU & 0.75$\pm$0.17 & 0.88$\pm$0.16 & 0.51$\pm$0.33 & 0.08 & 0.16 \\
\midrule
\multicolumn{6}{@{}l}{\textit{Fusion of standard encoders}} \\
EEGNet + ET-LSTM & 0.73$\pm$0.19 & 0.81$\pm$0.17 & 0.46$\pm$0.36 & 0.10 & 0.18 \\
ShallowConvNet + ET-LSTM & 0.72$\pm$0.19 & 0.83$\pm$0.17 & 0.43$\pm$0.36 & 0.08 & 0.17 \\
\midrule
\multicolumn{6}{@{}l}{\textit{Fusion architectures}} \\
Cross-Attention & 0.77$\pm$0.17 & 0.87$\pm$0.15 & 0.54$\pm$0.34 & 0.08 & 0.15 \\
DynamicGAT + ET Transformer & 0.77$\pm$0.19 & 0.88$\pm$0.16 & 0.54$\pm$0.37 & 0.07 & 0.16 \\
Multimodal Transformer & 0.77$\pm$0.16 & 0.84$\pm$0.18 & 0.53$\pm$0.31 & 0.15 & 0.18 \\
Dual Transformer & 0.76$\pm$0.19 & 0.86$\pm$0.17 & 0.54$\pm$0.37 & 0.08 & 0.16 \\
\textbf{Audited decoder (INS-HDGS-CMT)} & \textbf{0.75$\pm$0.18} & \textbf{0.88$\pm$0.17} & \textbf{0.49$\pm$0.36} & \textbf{0.14} & \textbf{0.18} \\
Audited decoder without the rule module & 0.74$\pm$0.17 & 0.87$\pm$0.18 & 0.48$\pm$0.34 & 0.17 & 0.18 \\
Early-fusion MLP & 0.73$\pm$0.19 & 0.85$\pm$0.19 & 0.46$\pm$0.38 & 0.06 & 0.17 \\
Late fusion (CNN-LSTM + ET-LSTM) & 0.73$\pm$0.20 & 0.83$\pm$0.18 & 0.45$\pm$0.39 & 0.12 & 0.19 \\
\botrule
\end{tabular*}
\end{table*}

\paragraph{The gaze inputs account for the accuracy.} Removing every gaze-derived input from the decoder, which leaves the gaze-free variant of audit step~2, lowers ROC-AUC by 0.287 (95\% CI $[0.207, 0.374]$; $p_{\mathrm{Holm}}<0.001$; $r_{\mathrm{rb}}=0.96$; 29 wins, 6 ties and 2 losses over folds), balanced accuracy by 0.195 ($[0.129, 0.261]$) and MCC by 0.391 (Supplementary Table~S14; Fig.~\ref{figladder}). The two gaze inputs are individually redundant: removing the gaze sequence while keeping the ROI saliency vector lowers ROC-AUC by 0.060 ($[-0.009, 0.138]$; $p_{\mathrm{Holm}}=0.12$), and removing the ROI vector while keeping the sequence by 0.015 ($p_{\mathrm{Holm}}=1.0$), so either gaze input alone carries most of what the label rewards, which is what its construction from five gaze statistics predicts. The per-participant picture behind these differences is given in Supplementary Fig.~S8. The gaze inputs therefore account for the decoder's accuracy on this index. What remains of the decoder without them, ROC-AUC 0.59, is the subject of the next question.

\begin{figure*}[!htbp]
\centering
\includegraphics[width=\textwidth]{fig_control_ladder_r2.pdf}
\caption{Every control against the audited decoder: mean paired difference over the 37 shared participants with its 95\% bootstrap CI (10{,}000 resamples); small points are individual participants, triangles those beyond the axis range. The raw two-sided Wilcoxon $p$ (zero differences discarded, as in Table~\ref{tab:abl}) is printed per row and filled markers denote $p<0.05$; Holm-corrected values are in Table~\ref{tab:abl} and Supplementary Table~S14. Balanced accuracy and MCC are at the transferred-threshold operating point.}
\label{figladder}
\end{figure*}

\subsection{No Tested EEG Encoder Outperformed a Probe on the Index's EEG Terms; the Recordings Carry Browsing State}\label{res:rq2}

\paragraph{No EEG encoder exceeded the linear probe on the EEG terms.} Table~\ref{tab:eeg} reports every encoder that receives no gaze input. The eight tuned encoders spanned balanced accuracy 0.48--0.59 and ROC-AUC 0.49--0.64: EEGNet reached 0.48 and 0.52, ShallowConvNet 0.59 and 0.62, DeepConvNet 0.59 and 0.63 and CNN-BiLSTM 0.57 and 0.64. The gaze-free variant of the decoder, dynamic graph attention with the spiking front-end and no gaze-derived input of any kind, reached balanced accuracy 0.55, ROC-AUC 0.59 and MCC 0.09, numerically below ShallowConvNet, DeepConvNet and CNN-BiLSTM and above EEGNet. None of the eight paired comparisons with the gaze-free variant was significant on any metric after Holm correction (balanced accuracy $p_{\mathrm{Holm}}\ge0.78$; MCC all 1.0; ROC-AUC $\ge0.44$; $|r_{\mathrm{rb}}|\le0.38$; no comparison exceeded 24 wins against 12 losses; Supplementary Table~S6). A Friedman test over the nine encoders rejected equality on balanced accuracy ($p=0.024$) and MCC ($p=0.030$) but not on ROC-AUC ($p=0.27$), an omnibus effect that does not by itself locate any pair; by average rank ShallowConvNet and DeepConvNet lead and EEGNet and GAT trail, and the Nemenyi diagram places all nine within one critical difference, in line with benchmark evidence for EEG decoding \cite{ref66} of each other (1.97 ranks; ShallowConvNet best at an average rank of 4.03, the gaze-free variant fourth at 4.97, EEGNet last at 5.92; Fig.~\ref{fig6}A), and the pooled curves of the gaze-free variant and the two best encoders lie close to the chance diagonal (Fig.~\ref{fig6}B). No encoder exceeded the 0.67 of the linear probe on the index's own EEG terms; the best, CNN-BiLSTM, reached 0.64. None of the tested EEG encoders has therefore been shown to use index-related information beyond the five terms the rule is built from; whether such information exists in the EEG is not decided by this comparison, which orders performance without identifying the information used. Two evaluation choices decide this result. An EEG branch that retains the gaze-derived ROI vector reaches balanced accuracy 0.70 and ROC-AUC 0.82 under the common budget and appears significantly better than seven of the eight encoders on balanced accuracy when those are trained under a fixed budget (paired Wilcoxon against the fixed-budget runs of Supplementary Table~S8, Holm-corrected over the eight; eight of eight on ROC-AUC); once the ROI input is removed and every baseline is tuned per fold, no advantage is detected in the corrected comparisons reported here.

\begin{table*}[!htbp]
\caption{EEG encoders under LOSOCV with no gaze input to any model (baselines tuned per fold); the gaze-free variant of the decoder is listed first. Mean $\pm$ SD over 37 folds at the transferred-threshold operating point; ECE and Brier pooled over the 347 held-out epochs. Common-budget values: Supplementary Table~S8.}\label{tab:eeg}
\footnotesize\setlength{\tabcolsep}{3pt}
\begin{tabular*}{\linewidth}{@{\extracolsep{\fill}}lccccccc@{}}
\toprule
Model & BalAcc & Macro-F1 & MCC & ROC-AUC & PR-AUC & ECE & Brier \\
\midrule
\textbf{Gaze-free variant of the decoder} & \textbf{0.55$\pm$0.17} & \textbf{0.51$\pm$0.29} & \textbf{0.09$\pm$0.33} & \textbf{0.59$\pm$0.26} & \textbf{0.66$\pm$0.26} & \textbf{0.05} & \textbf{0.25} \\
ShallowConvNet \cite{ref51} & 0.59$\pm$0.19 & 0.47$\pm$0.30 & 0.16$\pm$0.37 & 0.62$\pm$0.27 & 0.67$\pm$0.27 & 0.22 & 0.30 \\
DeepConvNet \cite{ref51} & 0.59$\pm$0.21 & 0.48$\pm$0.27 & 0.15$\pm$0.39 & 0.63$\pm$0.26 & 0.70$\pm$0.24 & 0.29 & 0.36 \\
CNN-BiLSTM & 0.57$\pm$0.15 & 0.40$\pm$0.32 & 0.13$\pm$0.32 & 0.64$\pm$0.30 & 0.70$\pm$0.27 & 0.20 & 0.32 \\
CNN-LSTM & 0.52$\pm$0.17 & 0.44$\pm$0.23 & 0.04$\pm$0.35 & 0.57$\pm$0.29 & 0.66$\pm$0.26 & 0.26 & 0.35 \\
EEG Transformer & 0.52$\pm$0.18 & 0.44$\pm$0.30 & 0.02$\pm$0.37 & 0.59$\pm$0.28 & 0.67$\pm$0.24 & 0.24 & 0.32 \\
TSception \cite{ref52} & 0.52$\pm$0.17 & 0.44$\pm$0.23 & 0.02$\pm$0.31 & 0.56$\pm$0.24 & 0.64$\pm$0.24 & 0.30 & 0.34 \\
GAT & 0.50$\pm$0.15 & 0.33$\pm$0.29 & 0.01$\pm$0.31 & 0.49$\pm$0.25 & 0.61$\pm$0.25 & 0.07 & 0.26 \\
EEGNet \cite{ref50} & 0.48$\pm$0.18 & 0.42$\pm$0.22 & $-0.04\pm$0.34 & 0.52$\pm$0.25 & 0.62$\pm$0.22 & 0.19 & 0.30 \\
\botrule
\end{tabular*}
\end{table*}

\begin{figure*}[!htbp]
\centering
\includegraphics[width=\textwidth]{fig6_combined.pdf}
\caption{Gaze-free EEG comparison, baselines tuned per fold. (A) Nemenyi critical-difference diagram over the nine EEG encoders (average rank on balanced accuracy, 37 folds; critical difference 1.97 at $\alpha=0.05$). (B) Pooled ROC and precision--recall curves of the gaze-free variant (mean per-fold ROC-AUC 0.59), CNN-BiLSTM (0.64) and DeepConvNet (0.63).}
\label{fig6}
\end{figure*}

\paragraph{No single marker carries the index either.} To test whether the index is carried by the canonical univariate EEG markers rather than by a multivariate pattern, each marker was evaluated against the labels with a within-participant label-permutation test (20{,}000 permutations over the 347 epochs; Fig.~\ref{fig7}A). No marker reached significance: the within-participant HIGH$-$LOW contrast was negligible for frontal-theta power (Cohen's $d=+0.09$, $p=0.18$), posterior-alpha power ($d=-0.03$, $p=0.40$) and fronto-posterior phase-locking in the theta ($d=+0.04$, $p=0.72$) and alpha ($d=-0.04$, $p=0.66$) bands, every effect lying within its permutation null and within the negligible range ($|d|<0.2$), so no within-participant difference is detected for any of the markers. The same epochs are separated at ROC-AUC 0.59 by the gaze-free variant and 0.88 by the decoder (Fig.~\ref{fig7}B). The contrast between the null markers and the gaze-free variant's 0.59 is consistent with index-related structure that is not concentrated in any single canonical marker, but it does not establish that such information is distributed across electrodes and bands, and it does not show that graph-structured integration is needed, since a linear encoder on the same band-power node features reaches 0.85 against the graph pathway's 0.86 on the positive control below and density-matched topology surrogates produce no detectable loss on this label (Section~\ref{res:rq3}).

\begin{figure*}[!htbp]
\centering
\includegraphics[width=\textwidth]{fig_concordance_depth.pdf}
\caption{No single marker carries the index. (A) Within-participant effect size (Cohen's $d$) of each canonical univariate marker with its permutation $p$ (20{,}000 permutations); grey band $|d|<0.2$. (B) ROC-AUC of the best single marker (fronto-posterior alpha PLV), the gaze-free variant and the audited decoder on the same 37 participants.}
\label{fig7}
\end{figure*}

\paragraph{The fold-wise label leaves the retested EEG results in place.} Because each held-out participant contributes to the pooled scaler and median of its own labels (Section~\ref{sec:label}), the principal runs were repeated under the fold-wise label with the matched validation rule (audit step~4; Supplementary Table~S17). The decoder and its gaze-free variant changed by amounts of the size of the run-to-run spread: the decoder gave balanced accuracy 0.715 against 0.747 (paired $\Delta=-0.033$, 95\% CI $[-0.092, +0.027]$, $p=0.32$), ROC-AUC 0.872 against 0.879 ($-0.007$, $[-0.039, +0.028]$, $p=0.47$) and MCC 0.455 against 0.485; the gaze-free variant gave 0.550, 0.605 and 0.091 against 0.552, 0.592 and 0.094 ($p\ge0.80$). The three reference encoders behave the same way. Under the fold-wise label EEGNet reaches balanced accuracy 0.53 and ROC-AUC 0.54 (pooled label 0.48 and 0.52; paired $\Delta$ $+0.051$ $[-0.014, +0.119]$ and $+0.019$ $[-0.081, +0.117]$, $p\ge0.21$), ShallowConvNet 0.60 and 0.64 (0.59 and 0.62; $+0.010$ and $+0.026$, $p\ge0.80$) and ET-LSTM 0.74 and 0.86 (0.76 and 0.89; $-0.018$ $[-0.065, +0.024]$, $p=0.53$, and $-0.031$ $[-0.057, -0.007]$, $p_{\mathrm{Holm}}=0.036$). The one change that exceeds the run-to-run spread is the gaze model's ranking, which falls by 0.03 in ROC-AUC once the held-out participant no longer contributes to the scaler and median of the gaze statistics from which the label is built; the model that sees the gaze terms carrying the label is the one the pooled construction had favoured (the rerun changes the label construction and the validation rule together, so the shift is not attributable to the label alone), and it remains 0.2 above every EEG encoder. The ordering, the gap between the EEG and gaze families and the position of every EEG encoder relative to the EEG-term probe are unchanged; the probes themselves, recomputed under the same fold-wise construction with the matched validation rule, give fold-mean ROC-AUC 0.68 on the five EEG terms and 0.93 on the five gaze terms, against 0.67 and 0.92 under the pooled label, and every fold remains two-class (Supplementary Table~S17). The values reported throughout remain those of the pooled label, under which the controls of Section~\ref{res:rq3} are paired; the fold-wise runs, which are also the runs in which validation participants are matched across models (Section~\ref{sec:eval}), show that no retested conclusion changes under the fold-wise construction.

\paragraph{The same recordings separate resting state from browsing.} The positive control (audit step~7) shows that the recordings and the pipelines carry condition-discriminating information (Fig.~\ref{figpos}; Supplementary Table~S16). On resting state versus brochure browsing, decoded gaze-free from the same participants under LOSOCV over 42 folds, tuned EEGNet and ShallowConvNet reach ROC-AUC 0.92 and 0.94, a linear encoder on the decoder's own band-power node features 0.85 (balanced accuracy 0.72), the graph pathway with a plain head 0.86, and the gaze-free variant with its MMD and adversarial regularisers retained and a seven-member ensemble 0.88 (balanced accuracy 0.72); the decoder with its gaze input reaches 0.99, since eye movement alone separates browsing from fixation. The spiking front-end alone reaches 0.55 in a configuration that differs from its production use (raw EEG pooled to 300 steps rather than band-power features over ten steps), and placed in front of the graph pathway with a plain head it lowers that pathway's 0.86 to 0.81; the gaze-free variant, which adds the subject-alignment regularisers and a larger ensemble, reaches 0.88 (ensemble sizes and regularisers per rung in Supplementary Table~S16). The EEG pathway of the decoder can therefore decode a contrast that is present in the recordings, a few points below the best convolutional encoder, and the same is true of every standard encoder. Placed side by side on the two labels (Fig.~\ref{figpos}), the encoders present on both, EEGNet, ShallowConvNet and the gaze-free variant, decode browsing state at 0.88--0.94 and the index at 0.52--0.62, at or below the EEG-term probe. The control does not isolate neural engagement: resting fixation and browsing differ in visual input, eye movements and muscle activity as well as in cortical state, so it establishes that usable condition-discriminating information is present in the recordings, and nothing further; the null on the index is not thereby shown to be a property of the label alone.

\begin{figure*}[!htbp]
\centering
\includegraphics[width=\textwidth]{fig_positive_control_r2.pdf}
\caption{The same EEG encoders on two labels from the same participants (Supplementary Table~S16); bars are mean per-fold ROC-AUC $\pm$ SD. (A) Positive control, resting state versus browsing (42 folds). (B) The engagement index (37 folds); red marks the only encoder that sees gaze (the audited decoder, 0.88). Blank positions are rungs not run on that track.}
\label{figpos}
\end{figure*}

\subsection{Architecture Case Study: What the Controls Show for Each Component}\label{res:rq3}

Each component was removed in turn and the model retrained on the same folds (audit step~5; Table~\ref{tab:abl}); Fig.~\ref{figladder} places every control on a common axis with bootstrap intervals so that effect sizes can be read against one another rather than from $p$-values alone. The gaze inputs carry the accuracy, as Section~\ref{res:rq1} showed ($-0.195$ in balanced accuracy and $-0.391$ in MCC when all are removed, $p<0.001$). Among the architectural components the picture is as follows.

\paragraph{The graph pathway's ablation cost is not attributable to the measured topology.} Removing the graph pathway produced the largest architectural degradation, $-0.092$ in balanced accuracy (raw $p=0.005$, $p_{\mathrm{Holm}}=0.041$) and $-0.169$ in MCC (raw $p=0.010$, $p_{\mathrm{Holm}}=0.077$). Whether this reflects the measured dynamic connectivity was tested with the two density-matched nulls, which leave node features, ROI vector and gaze untouched and replace only the adjacency: a static graph of the epoch's most persistent edges and, per window, a random graph of the same density. Neither produced a detectable change (balanced accuracy $-0.007$, $p=0.74$, and $+0.001$, $p=0.96$; MCC $+0.012$ and $+0.024$). Retrained on either surrogate, the pathway reaches the decoder's performance without the measured connectivity, so the ablation cost of 0.09 cannot be attributed to the time-varying topology the pathway was built to capture. The retrained surrogates may reach that performance through the processing the pathway applies to the per-window band-power features or through a different solution; the controls do not show what the original trained model uses, only that the measured topology was not necessary for the reported accuracy. This is consistent with Section~\ref{res:rq2}, where a linear encoder on the same node features reached 0.85 against the graph pathway's 0.86 on the positive control.

\paragraph{The spiking front-end is a representation with a measured cost.} Its removal changed balanced accuracy by $-0.046$ (raw $p=0.11$, $p_{\mathrm{Holm}}=0.52$), not significant after correction, so no claim of predictive necessity is made for it. The positive control shows what it does on its own: trained alone on the raw EEG epochs pooled to 300 LIF steps (42 folds) it separates resting state from browsing at ROC-AUC 0.55, against 0.92 and 0.94 for EEGNet and ShallowConvNet on the same recordings. That configuration differs from the production module, which receives band-power features over ten steps, so the 0.55 bounds what the encoder extracts from raw EEG on its own rather than what the production front-end discards; in front of the graph pathway, the three-member spiking-plus-graph configuration reaches 0.81 against 0.86 for the seven-member graph-only configuration, a comparison that differs in ensemble size and so does not isolate the spiking module's effect. The timed fold-1 checkpoint of the front-end fires at a mean rate of 7.5\%, a sparsity that would reduce energy only on event-driven neuromorphic hardware, on which it has not been run, so no hardware-energy figure is claimed or estimated. On the hardware used, every LIF step executes a dense matrix product, and direct timing shows the spiking encoder slower than a dense twin of identical widths: 2.70 against 0.45~ms per single-epoch inference on the CPU (3.45 against 0.69~ms at batch size 32) and 3.99 against 0.41~ms on an NVIDIA A100, on which neither encoder raised the board's power by more than 3~W above its 85~W idle draw, so the incremental inference energy could not be resolved (Supplementary Table~S9). The front-end is retained as an event-driven representation; an efficiency benefit on event-driven hardware was not evaluated in this study.

\paragraph{No detectable difference was found for the cross-modal transformer.} Removing the fusion transformer changed balanced accuracy by $-0.033$ ($p_{\mathrm{Holm}}=0.52$), and replacing it by a bidirectional EEG--eye-tracking cross-attention block without the graph token changed ROC-AUC by 0.012 in the variant's favour ($p_{\mathrm{Holm}}=1.0$; Supplementary Table~S14). Together with the comparisons of Table~\ref{tab:gaze}, in which the decoder differs from none of the fusion designs after correction, the two analyses detect no advantage of the fusion mechanism over simpler designs and no cost of removing it beyond the run-to-run spread, in a pipeline in which the gaze inputs already supply the discriminative structure.

\paragraph{The rule module's decisions are bypass-dominated.} Removing the module did not change accuracy ($-0.006$ in balanced accuracy, $p=0.59$), consistent with its role as an explanation layer, but the fidelity test (audit step~6; Supplementary Table~S13) shows that in the gated model the decision is dominated by the bypass. The learned gate never left its initialisation ($\alpha=0.570$, fold range 0.567--0.573), the rule activations stayed uniform ($a_r=0.125$ for all eight rules on HIGH and LOW epochs alike), the rule-only decision $\arg\max R$ agreed with the gated decision on 39\% of held-out epochs, and the rule margin was negatively correlated with the final logit margin ($r=-0.69$). Forcing $\alpha$ to zero at inference gave balanced accuracy 0.39 and ROC-AUC 0.22 against 0.72 and 0.88 for the gated decision at the calibrated operating point, whereas forcing it to one gave 0.725 and 0.877, with no detectable difference from the gated decision. With $\alpha=0.57$ the rule evidence still enters Eq.~\eqref{eq18} with weight 0.43, so this is bypass dominance rather than an algebraic absence of rule influence; the negative correlation of the rule margin with the decision shows that what the rules contribute is not aligned with it. A model trained throughout with the gate closed ($\alpha\equiv0$) reached balanced accuracy 0.715 and ROC-AUC 0.839 against 0.747 and 0.879 (Wilcoxon $p=0.047$ and $0.014$; at the calibrated operating point the balanced-accuracy difference is $+0.014$ in the variant's favour, $p=0.50$), so a rule-mediated decision is available at a cost of 0.04 in ROC-AUC, with a balanced-accuracy cost that depends on the operating point, and the attribution summaries of Section~\ref{res:rules} are drawn from that variant.

\paragraph{The remaining components are not significant after correction.} Removing the eye-tracking sequence branch while keeping the ROI vector ($-0.051$), the ROI gate and graph modulation ($-0.034$) or the MMD alignment term ($-0.031$) reduced balanced accuracy by amounts that are not significant after correction ($p_{\mathrm{Holm}}\ge0.25$) and lie between one and two times the run-to-run spread; removing the contrastive objective left it unchanged ($+0.005$). The sequence branch and the ROI vector are redundant with each other because either alone carries most of the gaze information the label rewards (Section~\ref{res:rq1}); the contrastive and MMD terms act on the representation geometry rather than on the decision boundary, and the MMD term is evaluated in few mini-batches (Section~\ref{sec:system}). The components implement the operations described in Section~\ref{sec:system}; their intended functions are distinct from the accuracy effects measured here, which are as stated.

\begin{table*}[!htbp]
\caption{Component ablation under LOSOCV (transferred-threshold operating point; decoder balanced accuracy $0.747$, MCC $0.485$), each row paired against the decoder. Paired two-sided Wilcoxon tests (zero differences discarded), Holm-adjusted within the balanced-accuracy and MCC families of the eight original ablations; the four rows outside that family carry raw $p$. ``$-$ Eye-tracking sequence branch'' retains the ROI saliency vector; ``$-$ All gaze-derived input'' is the gaze-free variant of Table~\ref{tab:eeg}; the graph-null rows replace only the adjacency by a density-matched static or random graph. $^{*}$Adjusted $p<0.05$.}\label{tab:abl}
\small
\begin{tabular*}{\linewidth}{@{\extracolsep{\fill}}lcccccc@{}}
\toprule
& \multicolumn{3}{c}{Balanced accuracy} & \multicolumn{3}{c}{MCC} \\
Configuration & $\Delta$ & $p_{\mathrm{raw}}$ & $p_{\mathrm{Holm}}$ & $\Delta$ & $p_{\mathrm{raw}}$ & $p_{\mathrm{Holm}}$ \\
\midrule
Audited decoder & --- & --- & --- & --- & --- & --- \\
$-$ Dynamic graph & $-0.092$ & 0.005 & \textbf{0.041}$^{*}$ & $-0.169$ & 0.010 & \textbf{0.077} \\
$-$ Eye-tracking sequence branch (ROI retained) & $-0.051$ & 0.035 & 0.25 & $-0.088$ & 0.058 & 0.33 \\
$-$ Spiking encoder & $-0.046$ & 0.111 & 0.52 & $-0.081$ & 0.105 & 0.42 \\
$-$ ROI gate and graph modulation & $-0.034$ & 0.068 & 0.41 & $-0.059$ & 0.044 & 0.31 \\
$-$ Fusion transformer & $-0.033$ & 0.103 & 0.52 & $-0.074$ & 0.056 & 0.33 \\
$-$ MMD alignment & $-0.031$ & 0.197 & 0.59 & $-0.047$ & 0.227 & 0.68 \\
$-$ Neuro-symbolic module & $-0.006$ & 0.586 & 1.00 & $-0.002$ & 0.850 & 1.00 \\
$-$ Contrastive objective & $+0.005$ & 0.887 & 1.00 & $+0.011$ & 0.906 & 1.00 \\
\midrule
$-$ All gaze-derived input (gaze-free variant) & $-0.195$ & $<0.001$ & -- & $-0.391$ & $<0.001$ & -- \\
Rule gate closed ($\alpha\equiv0$, rule-only decision) & $-0.033$ & 0.047 & -- & $-0.048$ & 0.093 & -- \\
Measured graph $\to$ density-matched static graph & $-0.007$ & 0.736 & -- & $+0.012$ & 0.959 & -- \\
Measured graph $\to$ density-matched random graph & $+0.001$ & 0.959 & -- & $+0.024$ & 0.959 & -- \\
\botrule
\end{tabular*}
\end{table*}

\subsection{Session-Level Dwell Predicts Product Selection; Nothing Tested Inside the Epoch Adds to It}\label{res:rq4}

The product-selection label of audit step~8 asks whether anything inside a 5-s epoch predicts whether the viewed product was selected for purchase (Table~\ref{tab:purchase}). Session-level dwell does: a logistic probe on the total time the product was looked at reaches a pooled ROC-AUC of 0.85 and a mean per-fold ROC-AUC of 0.88 over the 42 participants, and adding the number of gaze runs, the number of views and the anchor run raises these to 0.90 and 0.91. Nothing inside the epoch approaches this. The decoder, trained on the product-selection label in its production configuration, reaches balanced accuracy 0.566, ROC-AUC 0.606 and MCC 0.112, numerically above a linear probe on the five gaze statistics of the engagement rule computed within the epoch (pooled 0.57, fold mean 0.57; paired fold difference $+0.037$, 95\% CI $[-0.004, +0.081]$, $p=0.095$, not significant) and above the probe on its five EEG statistics (pooled 0.51, fold mean 0.53; $+0.080$, $[+0.023, +0.136]$, $p=0.012$), reference features that are not terms of this label; the five standard models trained on the same epochs reach ROC-AUC 0.485 (EEGNet), 0.488 (ShallowConvNet), 0.519 (ET-LSTM), 0.482 (EEGNet+ET-LSTM) and 0.524 (ShallowConvNet+ET-LSTM), with balanced accuracy 0.48--0.51 and MCC between $-0.03$ and $0.00$, and the two convolutional encoders show no detectable difference from each other, alone or fused ($p_{\mathrm{Holm}}\ge0.26$ on every metric). The decoder's margin over all five is significant (ROC-AUC $+0.082$ to $+0.125$, 95\% CIs excluding zero, $p_{\mathrm{Holm}}\le0.005$; balanced accuracy $+0.060$ to $+0.081$, $p_{\mathrm{Holm}}\le0.004$; MCC $+0.109$ to $+0.144$, $p_{\mathrm{Holm}}\le0.003$; Holm correction over the five comparisons within each metric), so it extracts within-epoch information that encoders trained from the raw sequences do not; the comparison does not identify which information. An incremental-validity probe asks whether that information adds to dwell. A LOSOCV logistic regression on log dwell alone reaches a mean per-fold ROC-AUC of 0.880 (pooled 0.854); the decoder's held-out epoch probability alone reaches 0.606 (0.591); the two together reach 0.885 (0.856), a gain of 0.005 (95\% CI $[-0.001, +0.012]$; Wilcoxon $p=0.11$; 22 folds up, 5 tied, 15 down); and the five standard models change the dwell probe by $-0.002$ to $+0.001$ (Table~\ref{tab:purchase}). The second-level regression uses held-out probabilities whose training folds included the other participants of each fold, a dependence that is expected, if anything, to favour the model; the stacking is not nested in the outer participant split, so the increment is not a leakage-free estimate. Two features of the design bound the reading. Total dwell is accumulated over the whole session, including viewing after the selection click, and the epoch anchor is chosen retrospectively as the longest gaze run on the page, so the session-level result is an association between looking time and selection rather than a prospective prediction from a common cutoff; the epoch content itself precedes the click by at least 250~ms. Within those bounds, what predicts product selection on these data is how long a product is looked at, a session-level behaviour by which the NeuMa product segments are themselves defined, and no tested model of a 5-s epoch adds detectably to it. The variants that remove gaze, connectivity or the ROI vector were not run on this label, so the source of the decoder's margin over the within-epoch reference probes on this label remains untested.

\begin{table*}[!htbp]
\caption{Product-selection label: whether the viewed product was selected for purchase (NeuMa item Q77), one 5-s epoch per participant--product pair anchored 1~s before the onset of the product's longest gaze run (2{,}338 epochs, 418 selected, 42 participants, LOSOCV). Models are trained on the epochs; probes are logistic regressions on within-epoch reference features or on session-level gaze behaviour, reported as pooled / mean per-fold ROC-AUC. The within-epoch probes use the five EEG and five gaze statistics of the engagement rule as reference features; they are not terms of this label. Mean $\pm$ SD over folds at the transferred-threshold operating point.}\label{tab:purchase}
\small
\begin{tabular*}{\linewidth}{@{\extracolsep{\fill}}p{0.48\linewidth}ccc@{}}
\toprule
Model or probe & BalAcc & ROC-AUC & MCC \\
\midrule
\multicolumn{4}{@{}l}{\textit{Models trained on the 5-s epoch}} \\
Audited decoder, production configuration & 0.566 $\pm$ 0.091 & 0.606 $\pm$ 0.118 & 0.112 \\
EEGNet & 0.506 $\pm$ 0.099 & 0.485 $\pm$ 0.120 & 0.003 \\
ShallowConvNet & 0.485 $\pm$ 0.097 & 0.488 $\pm$ 0.136 & $-0.032$ \\
ET-LSTM & 0.501 $\pm$ 0.075 & 0.519 $\pm$ 0.108 & $-0.001$ \\
EEGNet + ET-LSTM & 0.495 $\pm$ 0.075 & 0.482 $\pm$ 0.128 & $-0.005$ \\
ShallowConvNet + ET-LSTM & 0.502 $\pm$ 0.080 & 0.524 $\pm$ 0.125 & $-0.002$ \\
\midrule
\multicolumn{4}{@{}l}{\textit{Linear probes, within the epoch}} \\
Five EEG statistics of the engagement rule & --- & 0.51 / 0.53 & --- \\
Five gaze statistics of the engagement rule & --- & 0.57 / 0.57 & --- \\
\midrule
\multicolumn{4}{@{}l}{\textit{Linear probes, session level}} \\
Total dwell on the product & --- & 0.85 / 0.88 & --- \\
Dwell + number of runs + views + anchor run & --- & 0.90 / 0.91 & --- \\
\midrule
\multicolumn{4}{@{}l}{\textit{Incremental validity: log dwell plus the model's held-out epoch probability (LOSOCV logistic regression)}} \\
Log dwell alone & --- & 0.854 / 0.880 & --- \\
Log dwell + decoder probability (gain $+0.005$, 95\% CI $[-0.001, +0.012]$) & --- & 0.856 / 0.885 & --- \\
Log dwell + probability of any of the five standard models & --- & 0.852--0.853 / 0.879--0.881 & --- \\
\botrule
\end{tabular*}
\end{table*}

\subsection{Architecture Case Study: Attribution Summaries of the Rule-Only Variant}\label{res:rules}

The summaries below are taken from the variant trained with the gate closed, whose decision passes through the rule module (Section~\ref{res:rq3}); the traces of the gated production model are bypass-dominated and are not reported as explanations. They are attribution summaries: integrated gradients attribute a rule's activation to the inputs relative to a zero baseline and do not yield logical conditions, so the IF--THEN form used below is a shorthand for the signed inputs to which each rule's activation is most sensitive, and its fidelity as a verbal rule has not been validated. Throughout, attributions are learned associations, not causal claims about neural mechanism. They are drawn from the decoder rather than from EEGNet because EEGNet decodes the index at ROC-AUC 0.52 (Table~\ref{tab:eeg}), so an attribution of its decisions on this label would explain a chance-level model; the EEG-level interpretation that depends on no encoder is the single-marker analysis of Section~\ref{res:rq2}, and attributions of EEGNet on the positive control, where it does carry signal, are the complement noted in Section~\ref{disc:prior}.

\paragraph{Attribution summaries.} Each rule's activation is attributed to the physical inputs by integrated gradients, giving a $19\times5$ electrode-by-band map per rule whose three largest terms, with their sign, form its attribution summary (Fig.~\ref{fig8}C,~D; all eight rules in Supplementary Fig.~S7). Across the eight rules, 47--85\% (mean 65\%) of the attribution mass falls on the EEG node features and the remainder on the gaze and ROI inputs; seven of the eight rules carry a horizontal-gaze term and one (rule 4) a vertical-gaze term, and the electrode--band terms that recur are posterior alpha at O1, O2 and Pz, with C3 gamma. Of the four LOW rules, which together dominate 64\% of epochs, three (rules 2, 3 and 5) are alpha-up and gaze-left; of the four HIGH rules, which dominate 36\%, three (rules 6, 7 and 8) are posterior-alpha-down and gaze-right (for example, rule 7: O1-alpha$\downarrow$ $\wedge$ O2-alpha$\downarrow$ $\wedge$ Pz-alpha$\downarrow$, with gaze-$x\uparrow$ $\rightarrow$ HIGH, read as a shorthand). Dominance is assigned by the largest of eight activations whose means are all close to $1/8$ (0.124--0.127 in the rule-only variant, Fig.~\ref{fig8}C; 0.125 in the gated model), so the rules partition the epochs by small activation margins rather than by sharply selective premises, and with near-uniform activations the module is close to an average of eight linear heads in the key space. The summaries therefore describe what each head's activation is sensitive to; they are not validated rules, the learned premises remain vectors in the key space, and rule stability across folds has not been quantified.

\begin{figure*}[!htbp]
\centering
\includegraphics[width=\textwidth]{fig5_explainability.pdf}
\caption{Explainability panels. (A) Mean graph attention received per electrode and (B) the electrode-by-electrode attention matrix of the production model for held-out participant S01 (fold-1 checkpoint); both are descriptive. (C) Mean soft-rule activations and (D) attribution summaries (three largest electrode--band terms, dominant eye-tracking stream, rule conclusion) of the variant trained with the rule gate closed, aggregated over the 347 held-out epochs of all 37 fold checkpoints with rules pooled by index, a correspondence across checkpoints that has not been validated; full maps in Supplementary Fig.~S7.}
\label{fig8}
\end{figure*}

\paragraph{Population-level attributions.} Integrated gradients aggregated across folds rank the ROI saliency vector first (normalised importance 0.52), posterior alpha power (0.18) and frontal theta power (0.16) as the leading EEG inputs, and gaze position (0.13) as the leading raw eye-tracking feature, with pupil dynamics (0.01) last; the connectivity input enters as a binary mask that receives no gradient, so integrated gradients do not assess it (Supplementary Table~S2). In the illustrated fold-1 checkpoint for S01, the mean received graph attention is largest at O2, F7 and the fronto-central electrodes F3 and C3, and smallest at the parietal electrodes Pz and P4 (Fig.~\ref{fig8}A,~B), a single-checkpoint description; because density-matched adjacencies reproduce the decoder's accuracy, these coefficients describe where the trained model attends, not what its decision depends on. Read with the single-marker result of Section~\ref{res:rq2}, the attributions are consistent with an index-related pattern that is not concentrated in a single marker; they do not establish how it is distributed, and the topology surrogates show only that the measured connectivity was not necessary for the reported accuracy.

\paragraph{The gaze phenotype.} Across the 385 epochs, HIGH epochs are set apart by broader, more exploratory gaze, with greater spatial dispersion (Cohen's $d=1.02$), longer scan-paths and faster gaze kinematics (both $d=0.44$), and modestly larger pupil dilation ($d=0.30$); LOW epochs show gaze confined to a narrow region (Fig.~\ref{fig9}). Because the label is built in part from gaze statistics, these descriptors characterise the behavioural signature of the label rather than predicting it independently.

\begin{figure*}[!htbp]
\centering
\includegraphics[width=\textwidth]{fig_et_phenotype.pdf}
\caption{Eye-tracking phenotype (385 epochs). (A,~B) Gaze-position density for HIGH and LOW epochs. (C) Pupil-diameter time course (mean $\pm$ 95\% CI). (D) Cohen's $d$ of the eye-tracking descriptors. The label is built partly from gaze statistics, so these descriptors characterise rather than predict engagement.}
\label{fig9}
\end{figure*}

\paragraph{Two cases.} Two held-out epochs, the most confident correctly classified epoch of each class for its participant, are traced through the same trained ensemble (Fig.~\ref{fig10}; regional EEG signatures in Supplementary Fig.~S2, ROI vectors side by side in Supplementary Fig.~S3). For the HIGH exemplar (S24), frontal-theta power is higher in the HIGH than in the LOW epochs of the same participant, the ROI vector is spread over eight of the ten grid cells across both rows, the gaze ranges widely over the page, and the attribution is distributed across modalities. For the LOW exemplar (S30), the within-participant theta contrast has the same sign while posterior alpha runs in the opposite direction to S24, echoing the population-level null of each marker; the defining feature is spatial, with the ROI vector concentrated in four cells of the upper row (62\% in one cell) and zero elsewhere, and the attribution led by the ROI trajectory. The two cases portray the HIGH-index exemplar as broad, multimodally coherent exploration of the page and the LOW-index exemplar as gaze confined to one region, in keeping with the gaze-dispersion term of the label, with frontal theta a case-level EEG correlate that is null at the population level.

\begin{figure*}[!htbp]
\centering
\includegraphics[width=\textwidth]{fig_gaze_pred.pdf}
\caption{HIGH (S24) and LOW (S30) exemplars: gaze trajectory, ROI saliency vector over the ten $5\times2$ grid cells (1--5 upper row, 6--10 lower row), predicted class probability, and integrated-gradients attribution over the three input streams that receive gradient (the connectivity input enters only as a binary mask and receives none).}
\label{fig10}
\end{figure*}

\subsection{Robustness Checks}\label{res:robust}

\paragraph{Participant-level behaviour and label-free thresholds.} Most held-out participants are decoded well above chance and a minority markedly lower (Supplementary Fig.~S1). A per-participant diagnostic that decomposes each failure into ranking, threshold and calibration components separates two modes. The larger group is threshold-limited but representationally sound: for these participants the decoder ranks the held-out epochs perfectly (ROC-AUC 1.00) yet scores near chance at the transferred threshold, so the discriminative information is present and the error is one of operating-point transfer. A smaller group reflects genuine representation failure. Consistent with the first mode, the decoder's ROC-AUC of 0.88 exceeds what its transferred-threshold balanced accuracy of 0.75 would imply. A label-free, prevalence-matched threshold per held-out participant, the quantile of the participant's own unlabelled predicted probabilities at the training-set class prior, gives balanced accuracy 0.740 and MCC 0.43 when computed transductively over all of a participant's epochs, and a strictly causal variant, in which the threshold for epoch $k$ is computed from the preceding epochs only after a short unlabelled warm-up at the transferred threshold, gives 0.748 and 0.46 with a warm-up of three epochs (15~s of analysed data) and 0.746 and 0.47 with five (Supplementary Table~S12). No test label enters either rule, so the operating point can be recovered prospectively, online, from a new participant's first unlabelled epochs. Retraining the decoder with per-epoch instance normalisation, which needs no participant-level statistic, gave balanced accuracy 0.763 against 0.747 (Wilcoxon $p=0.48$), so no model-side step depends on statistics of the new participant's recording; the preprocessing of Section~\ref{sec:preproc}, filtering, independent component analysis and gap interpolation, is fitted on each whole recording offline, and a prospective deployment would require causal versions of those steps, which were not evaluated. Three balanced-accuracy operating points thus describe the same ranking: 0.75 at the transferred threshold, 0.72 at the calibrated point and 0.74--0.75 under the label-free rules.

\paragraph{Label regime and participant coverage.} Splitting each participant at their own rank-based median renders all 42 participants evaluable and poses a within-participant discrimination task; under it the ROI-retained EEG branch (the EEG pathway with the ROI saliency vector, common budget) gave balanced accuracy 0.69 over 42 participants against 0.70 over 37 and ROC-AUC 0.80 against 0.82 under the common budget (Supplementary Table~S1), so the coverage regime gives similar aggregate performance; this does not establish that the comparisons under the headline label are free of selection effects, and their estimand is performance over participants whose epochs contain both classes. The check bears on coverage only; no conclusion about gaze-free decoding rests on it. The fold-wise relabelling is reported in Section~\ref{res:rq2}.

\paragraph{Connectivity measure and threshold.} Under phase-locking value in place of Pearson correlation (broadband and five band-limited variants), evaluated on a reduced subset of 16 participants, all seven measures gave near-chance discrimination of the index (mean per-participant ROC-AUC 0.47--0.60; Supplementary Table~S3), as the gaze-free variant does at full scale under Pearson correlation; the comparison cannot favour a measure, and the topology surrogates of Section~\ref{res:rq3} bear on this: on this label the measured connectivity was not necessary for the decoder's accuracy, whichever measure produces it. The connectivity threshold $\tau$ acts as a smooth density control (Supplementary Table~S7, Fig.~S4): over the 3{,}850 windows the retained fraction of the 171 electrode pairs falls from 0.85 at $\tau=0.10$ through 0.57 at the value used (mean degree 10.3) to 0.15 at $\tau=0.70$, while fragmentation, electrodes left with fewer than three edges, stays below one per window up to $\tau=0.30$ and rises steeply beyond (2.0 at 0.40, 4.9 at 0.50, 13.2 at 0.70), which is the structural rationale for the a priori value. Retraining the gaze-free variant for $\tau\in\{0.20,\dots,0.50\}$ under a reduced budget gave balanced accuracy 0.54--0.59 and ROC-AUC 0.52--0.59, flat across the range; the threshold is structurally motivated and predictively immaterial.

\paragraph{ROI grid and cohort size.} Rebuilding the ROI saliency vector on grids from $2\times1$ to $10\times8$ cells changes its granularity (mean largest-cell share 0.74 to 0.23) but preserves the per-epoch ordering of attentional concentration (Spearman correlation with the $5\times2$ vector 0.71--0.82 for every grid of six or more cells, 0.56 for the two-cell grid); the label contains no spatial grid, so it is unaffected; and the decoder retrained on $2\times1$, $3\times2$, $6\times4$ and $8\times6$ grids gave balanced accuracy 0.72--0.76 and ROC-AUC 0.88--0.90, none differing from the $5\times2$ run (smallest $p=0.53$; Supplementary Table~S11, Fig.~S6). Repeating LOSOCV on random subsets of 10, 16, 24, 30 and 37 participants (three draws each) gave held-out ROC-AUC of 0.86 at 10 participants, 0.87 at 16 and 0.88 at 37, with balanced accuracy rising from 0.68 to 0.74 and a slope between 30 and 37 participants of 0.001 in ROC-AUC and 0.002 in balanced accuracy per additional participant (Supplementary Fig.~S5). Ranking quality has plateaued by 16 participants within the observed range; a subsampling curve cannot say what a larger cohort would give, but the small gain per participant is consistent with performance limited by the label's construction rather than by the amount of training data, which the curve cannot establish.
